% !TEX root = ../[topic]-notes.tex
% Chapter Outline Template - ~30% filled

%======================================================================
% \chapter{章节标题}\label{ch:X}
%======================================================================

\begin{refsection}
\chapter{第5章 协变量平衡与因果回归}\label{ch:5}

%======================================================================
\section{引言}\label{sec:5-intro}
%====================================================================--

在第三章中，我们学习了潜在结果框架的基本概念%
\addtocounter{footnote}{1}\addtocounter{footnote}{-1}\footnotesize%
\par\_notes{关于潜在结果的详细讨论见 \cref{sec:3-1}。}，
认识到因果效应定义为%
$\tau = \mathbb{E}[Y(1) - Y(0)]$。

然而，这一框架在实践中面临一个根本性困难：\textbf{无法同时观测每个单元的两种潜在结果}。
这引导我们思考：如何在只有部分观测数据的情况下，估计因果效应？

\textbf{本章结构概览}：
\begin{enumerate}
\item \S\ref{sec:5-1} 引入协变量平衡的核心概念；
\item \S\ref{sec:5-2} 建立协变量平衡与因果识别的主要定理；
\item \S\ref{sec:5-3} 通过数值例子说明定理的应用；
\item \S\ref{sec:5-4} 讨论协变量平衡的诊断方法。
\end{enumerate}

%======================================================================
\section{协变量平衡的定义}\label{sec:5-1}
%====================================================================--

\subsection{从混杂因子谈起}\label{sec:5-1-1}

\textbf{核心问题}：当处理分配 $Z$ 与潜在结果 $(Y(1), Y(0))$ 相关时，
即存在\textbf{混杂因子}（confounder）时，直接比较处理组和对照组会产生偏误。

\textbf{历史脉络}：这个问题最早由...


\subsection{形式化定义}\label{sec:5-1-2}

在正式定义协变量平衡之前，我们考虑一个简单例子说明其直观含义...

\begin{Example}[二元处理的协变量平衡]\label{ex:5-1}
设处理变量 $Z \in \{0, 1\}$，协变量 $X \in \mathbb{R}^p$...
\end{Example}
\par\_notes{完整计算过程见附录 \cref{sec:appendix-ex-5-1}。}

\begin{Definition}[5.1 — 协变量平衡]\label{def:5-1}
我们称处理组和对照组满足\textbf{协变量平衡}，如果
\begin{equation}\label{eq:5-1-balance}
\mathbb{P}(Z=1 \mid X) = \mathbb{P}(Z=1)
\end{equation}
或者更一般地，如果
\begin{equation}\label{eq:5-1-balance-general}
\mathbb{E}[Z \mid X] = \mathbb{E}[Z]。
\end{equation}
\end{Definition}

\textbf{直观理解}：协变量平衡意味着，处理分配的倾向与协变量无关，
从而使我们可以将组间差异归因于处理效应而非混淆因素。


%======================================================================
\section{协变量平衡与因果识别}\label{sec:5-2}
%====================================================================--

\subsection{主要定理}\label{sec:5-2-1}

\begin{Theorem}[5.1 — 协变量平衡与因果识别]\label/thm:5-1}
假设以下条件成立：
\begin{enumerate}
\item \textbf{强可忽略性}（Strong Ignorability）：
$(Y(1), Y(0)) \Perp Z \mid X$；
\item \textbf{重叠性}（Overlap）：$0 < \mathbb{P}(Z=1 \mid X) < 1$。
\end{enumerate}
则有
\begin{equation}\label{eq:5-2-ate}
\mathbb{E}[Y(1) - Y(0)] = \mathbb{E}_X\big[\mathbb{E}(Y \mid Z=1, X) - \mathbb{E}(Y \mid Z=0, X)\big]。
\end{equation}
\end{Theorem}

\begin{Proof}
证明概要：证明的核心思想是...

\done
\par\_notes{完整证明见附录 \cref{sec:appendix-thm-5-1}。}
\end{Proof}


%======================================================================
\section{数值例子}\label{sec:5-3}
%====================================================================--

\begin{Example}[协变量平衡的计算]\label{ex:5-2}
考虑一个简单研究...
\end{Example}
\par\_notes{完整计算过程见附录 \cref{sec:appendix-ex-5-2}。}


%======================================================================
\section{协变量平衡的诊断}\label{sec:5-4}
%====================================================================--

\textbf{动机}：在实际研究中，我们通常无法验证协变量平衡是否成立，
但可以通过一些诊断方法来评估...


%======================================================================
\section{小结}\label/sec:5-summary}
%====================================================================--

本章引入了协变量平衡的核心概念...

\textbf{主要结果}：
\begin{enumerate}
\item 定义 \ref{def:5-1} 提供了协变量平衡的数学刻画；
\item 定理 \ref/thm:5-1} 表明，在强可忽略性和重叠性下...
\end{enumerate}

在下一章中，我们将学习如何通过...来评估和实现协变量平衡。

\par\_notes{主要定理的完整证明见附录 \crefrange{sec:appendix-thm-5-1}{sec:appendix-thm-5-2}。}

\end{refsection}
\endinput
